% =====================================================================
\section{Conclusion}\label{sec:conclusion}
Using \num{626679} forecast cases from 517 GDP drifters in the Gulf Stream, we evaluated two physical baselines and ten learned residual-correction models on 78 held-out drifters. Recent motion history substantially improves 24-hour position forecasts. History-enhanced Gradient Boosting reduced the median error of Constant Velocity from \SI{15.9}{km} to \SI{9.2}{km}, and history improved every model family, most strongly the nonlinear tree ensembles. Forecast difficulty is governed mainly by recent speed. Turning is associated with lower error because turning drifters tend to be slow, and the benefit of history is positive in every motion regime and largest for fast, turning drifters. The advantage is geographically widespread and stable across 2007--2022, and adding SST as a predictor did not help. Practically, the results show that a drifter's own recent track substantially narrows where it is likely to be a day later, and that recent speed indicates in advance which forecasts deserve least confidence. These conclusions concern retrospective forecasts from smoothed hourly data that mix drogued and undrogued drifters. Real-time inputs, drogue-aware evaluation and multiple lead times are the essential next steps toward an operational assessment.
